
% abstract.tex — Full proposal draft, Week 8
% Place this file alongside main_draft.tex and 
% abstract.tex — Full proposal draft, Week 8
% Place this file alongside main_draft.tex and 
% abstract.tex — Full proposal draft, Week 8
% Place this file alongside main_draft.tex and 
% abstract.tex — Full proposal draft, Week 8
% Place this file alongside main_draft.tex and \input{abstract} where needed.

\begin{abstract}

Individual differences in sustained attention predict academic achievement, occupational performance, and vulnerability to clinical conditions, yet the neuroimaging tools best suited to studying these differences are expensive, immobile, and concentrated in well-funded institutions. This project asks whether a brief, portable brain recording session can approximate whole-brain neural measurement without a scanner. Specifically, I test whether functional near-infrared spectroscopy (fNIRS) signals recorded during naturalistic movie-watching can train a computational model that predicts whole-brain fMRI activity, and whether that model generalizes to structured cognitive tasks where it can detect meaningful individual differences in attentional control.

The approach uses adapted principal component regression (aPCR) to learn a linear mapping between 42 fNIRS channels and 122 fMRI brain regions, trained on temporally aligned data from two separate participant groups (N\,=\,28 fNIRS, N\,=\,59 fMRI) who watched the same film. The trained model is then frozen and applied to fNIRS data from cognitive tasks (n-back, Continuous Performance Task) that the model has never encountered. The primary outcome is whether predicted whole-brain activation patterns correlate with behavioral measures of attentional stability, operationalized as intraindividual reaction time variability.

Preliminary results from the movie-watching training phase show that the model significantly predicts 56 of 122 brain regions (FDR $q < 0.05$), with particularly strong performance in the dorsal attention network (11 of 14 ROIs significant) and ventral attention network (14 of 24 ROIs). The predicted inter-subject functional connectivity matrix closely matches the observed fMRI connectivity structure ($r = 0.617$, $p = 0.001$). These results establish the model's capacity to recover large-scale brain organization from surface signals alone, providing the foundation for the cross-task individual differences analyses that follow.

\end{abstract}
 where needed.

\begin{abstract}

Individual differences in sustained attention predict academic achievement, occupational performance, and vulnerability to clinical conditions, yet the neuroimaging tools best suited to studying these differences are expensive, immobile, and concentrated in well-funded institutions. This project asks whether a brief, portable brain recording session can approximate whole-brain neural measurement without a scanner. Specifically, I test whether functional near-infrared spectroscopy (fNIRS) signals recorded during naturalistic movie-watching can train a computational model that predicts whole-brain fMRI activity, and whether that model generalizes to structured cognitive tasks where it can detect meaningful individual differences in attentional control.

The approach uses adapted principal component regression (aPCR) to learn a linear mapping between 42 fNIRS channels and 122 fMRI brain regions, trained on temporally aligned data from two separate participant groups (N\,=\,28 fNIRS, N\,=\,59 fMRI) who watched the same film. The trained model is then frozen and applied to fNIRS data from cognitive tasks (n-back, Continuous Performance Task) that the model has never encountered. The primary outcome is whether predicted whole-brain activation patterns correlate with behavioral measures of attentional stability, operationalized as intraindividual reaction time variability.

Preliminary results from the movie-watching training phase show that the model significantly predicts 56 of 122 brain regions (FDR $q < 0.05$), with particularly strong performance in the dorsal attention network (11 of 14 ROIs significant) and ventral attention network (14 of 24 ROIs). The predicted inter-subject functional connectivity matrix closely matches the observed fMRI connectivity structure ($r = 0.617$, $p = 0.001$). These results establish the model's capacity to recover large-scale brain organization from surface signals alone, providing the foundation for the cross-task individual differences analyses that follow.

\end{abstract}
 where needed.

\begin{abstract}

Individual differences in sustained attention predict academic achievement, occupational performance, and vulnerability to clinical conditions, yet the neuroimaging tools best suited to studying these differences are expensive, immobile, and concentrated in well-funded institutions. This project asks whether a brief, portable brain recording session can approximate whole-brain neural measurement without a scanner. Specifically, I test whether functional near-infrared spectroscopy (fNIRS) signals recorded during naturalistic movie-watching can train a computational model that predicts whole-brain fMRI activity, and whether that model generalizes to structured cognitive tasks where it can detect meaningful individual differences in attentional control.

The approach uses adapted principal component regression (aPCR) to learn a linear mapping between 42 fNIRS channels and 122 fMRI brain regions, trained on temporally aligned data from two separate participant groups (N\,=\,28 fNIRS, N\,=\,59 fMRI) who watched the same film. The trained model is then frozen and applied to fNIRS data from cognitive tasks (n-back, Continuous Performance Task) that the model has never encountered. The primary outcome is whether predicted whole-brain activation patterns correlate with behavioral measures of attentional stability, operationalized as intraindividual reaction time variability.

Preliminary results from the movie-watching training phase show that the model significantly predicts 56 of 122 brain regions (FDR $q < 0.05$), with particularly strong performance in the dorsal attention network (11 of 14 ROIs significant) and ventral attention network (14 of 24 ROIs). The predicted inter-subject functional connectivity matrix closely matches the observed fMRI connectivity structure ($r = 0.617$, $p = 0.001$). These results establish the model's capacity to recover large-scale brain organization from surface signals alone, providing the foundation for the cross-task individual differences analyses that follow.

\end{abstract}
 where needed.

\begin{abstract}

Individual differences in sustained attention predict academic achievement, occupational performance, and vulnerability to clinical conditions, yet the neuroimaging tools best suited to studying these differences are expensive, immobile, and concentrated in well-funded institutions. This project asks whether a brief, portable brain recording session can approximate whole-brain neural measurement without a scanner. Specifically, I test whether functional near-infrared spectroscopy (fNIRS) signals recorded during naturalistic movie-watching can train a computational model that predicts whole-brain fMRI activity, and whether that model generalizes to structured cognitive tasks where it can detect meaningful individual differences in attentional control.

The approach uses adapted principal component regression (aPCR) to learn a linear mapping between 42 fNIRS channels and 122 fMRI brain regions, trained on temporally aligned data from two separate participant groups (N\,=\,28 fNIRS, N\,=\,59 fMRI) who watched the same film. The trained model is then frozen and applied to fNIRS data from cognitive tasks (n-back, Continuous Performance Task) that the model has never encountered. The primary outcome is whether predicted whole-brain activation patterns correlate with behavioral measures of attentional stability, operationalized as intraindividual reaction time variability.

Preliminary results from the movie-watching training phase show that the model significantly predicts 56 of 122 brain regions (FDR $q < 0.05$), with particularly strong performance in the dorsal attention network (11 of 14 ROIs significant) and ventral attention network (14 of 24 ROIs). The predicted inter-subject functional connectivity matrix closely matches the observed fMRI connectivity structure ($r = 0.617$, $p = 0.001$). These results establish the model's capacity to recover large-scale brain organization from surface signals alone, providing the foundation for the cross-task individual differences analyses that follow.

\end{abstract}
